\documentclass[border=15pt, multi, tikz]{standalone} 
\usepackage{import}
\subimport{../../layers/}{init}
\usetikzlibrary{positioning}
\usetikzlibrary{3d}

% --- GRAYSCALE / BLACK & WHITE PALETTE ---
\def\ConvColor{black!10}      
\def\PoolColor{black!30}      
\def\FcColor{black!55}        
\def\SoftmaxColor{black!80}   

\newcommand{\copymidarrow}{\tikz \draw[-Stealth,line width=0.8mm,draw=black] (-0.3,0) -- ++(0.3,0);}

\begin{document}
\begin{tikzpicture}
\tikzstyle{connection}=[ultra thick,every node/.style={sloped,allow upside down},draw=\edgecolor,opacity=0.7]

% --- TEXT FORMATTING (Reduced text width to prevent overlap) ---
\tikzstyle{labeltext}=[
    draw, 
    thin, 
    rounded corners=2pt, 
    inner sep=3pt,           % Reduced padding to save horizontal space
    align=center, 
    font=\normalsize, 
    anchor=north, 
    text width=2.8cm,        % REDUCED WIDTH: Forces text to wrap and makes boxes narrower
    yshift=-2.5cm
]

% --- DIAGRAM HEADING ---
\node[font=\Large\bfseries, align=center] at (17,4) {Deep Learning Convolutional Neural Network (CNN) Architecture};

% 1. Input Image
\pic[shift={(0,0,0)}] at (0,0,0) 
    {Box={
        name=input,
        caption= ,          % DO NOT put "Input Image" here; keep it empty
        xlabel={{3, }}, 
        zlabel=100, 
        fill=\ConvColor,
        height=30,
        width=2, 
        depth=30
        }
    };

% This node forces "Input Image" to stay on ONE line and sit above the rectangle
\node[anchor=south, font=\normalsize\bfseries, yshift=-2.4cm, xshift=-0.8cm] at (input-south) {Original Image};

% This rectangle contains the technical specs
\node[labeltext, xshift=-0.8cm] at (input-south) {RGB Image\\100$\times$100$\times$3};

% 2. BLOCK 1
\pic[shift={(3,0,0)}] at (input-east) 
    {Box={
        name=conv1,
        caption=Block 1,
        xlabel={{32, }},
        zlabel=100,
        fill=\ConvColor,
        height=30,
        width=3,
        depth=30
        }
    };
\pic[shift={(0,0,0)}] at (conv1-east) 
    {Box={
        name=pool1,
        caption= ,
        fill=\PoolColor,
        opacity=0.8,
        height=15,
        width=3,
        depth=15
        }
    };
\draw [connection]  (input-east) -- node {\midarrow} (conv1-west);
\node[labeltext, xshift=-0.8cm] at (conv1-south) {Conv2D (3x3)\\+ BN + ReLU\\+ MaxPool (2x2)};

% 3. BLOCK 2
\pic[shift={(3,0,0)}] at (pool1-east) 
    {Box={
        name=conv2,
        caption=Block 2,
        xlabel={{64, }},
        zlabel=50,
        fill=\ConvColor,
        height=15,
        width=5,
        depth=15
        }
    };
\pic[shift={(0,0,0)}] at (conv2-east) 
    {Box={
        name=pool2,
        caption= ,
        fill=\PoolColor,
        opacity=0.8,
        height=7.5,
        width=5,
        depth=7.5
        }
    };
\draw [connection]  (pool1-east) -- node {\midarrow} (conv2-west);
\node[labeltext, xshift=-0.8cm] at (conv2-south) {Conv2D (3x3)\\+ BN + ReLU\\+ MaxPool (2x2)};

% 4. BLOCK 3
\pic[shift={(3,0,0)}] at (pool2-east) 
    {Box={
        name=conv3,
        caption=Block 3,
        xlabel={{128, }},
        zlabel=25,
        fill=\ConvColor,
        height=7.5,
        width=8,
        depth=7.5
        }
    };
\pic[shift={(0,0,0)}] at (conv3-east) 
    {Box={
        name=pool3,
        caption= ,
        fill=\PoolColor,
        opacity=0.8,
        height=3.5,
        width=8,
        depth=3.5
        }
    };
\draw [connection]  (pool2-east) -- node {\midarrow} (conv3-west);
\node[labeltext] at (conv3-south) {Conv2D (3x3)\\+ BN + ReLU\\+ MaxPool (2x2)};

% 5. BLOCK 4
\pic[shift={(3,0,0)}] at (pool3-east) 
    {Box={
        name=conv4,
        caption=Block 4,
        xlabel={{256, }},
        zlabel=12,
        fill=\ConvColor,
        height=3.5,
        width=12,
        depth=3.5
        }
    };
\pic[shift={(0,0,0)}] at (conv4-east) 
    {Box={
        name=pool4,
        caption= ,
        fill=\PoolColor,
        opacity=0.8,
        height=1.5,
        width=12,
        depth=1.5
        }
    };
\draw [connection]  (pool3-east) -- node {\midarrow} (conv4-west);
\node[labeltext] at (conv4-south) {Conv2D (3x3)\\+ BN + ReLU\\+ MaxPool (2x2)};

% 6. DENSE 1 
\pic[shift={(3,0,0)}] at (pool4-east) 
    {Box={
        name=Fc1,
        caption=Dense 1,
        xlabel={{1, }},
        zlabel=128,
        fill=\FcColor,
        height=1.5,
        width=1.5,
        depth=15
        }
    };
\draw [connection]  (pool4-east) -- node {\midarrow} (Fc1-west);
\node[labeltext] at (Fc1-south) {\textbf{Flatten} - Dense\\+ L2 (0.001) + BN\\+ ReLU + Dropout(0.5)};

% 7. DENSE 2
\pic[shift={(4,0,0)}] at (Fc1-east) 
    {Box={
        name=Fc2,
        caption=Dense 2,
        xlabel={{1, }},
        zlabel=64,
        fill=\FcColor,
        height=1.5,
        width=1.5,
        depth=8
        }
    };
\draw [connection]  (Fc1-east) -- node {\midarrow} (Fc2-west);
\node[labeltext] at (Fc2-south) {Dense (64)\\+ L2 (0.001)\\+ BN + ReLU};

% 8. OUTPUT
\pic[shift={(3,0,0)}] at (Fc2-east) 
    {Box={
        name=soft1,
        caption=Output,
        xlabel={{" ","dummy"}},
        zlabel=8,
        fill=\SoftmaxColor,
        opacity=0.9,
        height=1.5,
        width=1.5,
        depth=3
        }
    };
\draw [connection]  (Fc2-east) -- node {\midarrow} (soft1-west);
\node[labeltext] at (soft1-south) {Softmax\\(8 Classes)};

\end{tikzpicture}
\end{document}